\documentclass[12pt]{ximera}
\title{Practice Exam 3}
\begin{document}
\begin{abstract}
Practice for Exam 3, covering Sections 8.1--8.5: properties of combinations and
expected value, choosing between permutations and combinations, Pascal's triangle,
counting supervillain teams, and building a binomial distribution to compute an
expected prize.
\end{abstract}
\maketitle

\section*{Practice Exam 3}

\begin{problem}
Assess whether each of the following claims is true or false.

\begin{prompt}
\textbf{(a)} Claim: for any positive whole number $n$, the value of $k$ that gives
the largest combination $C(n,k)$ is approximately half of $n$.

\begin{multipleChoice}
\choice[correct]{True}
\choice{False}
\end{multipleChoice}

\begin{feedback}
True. Each row of Pascal's triangle rises to a peak in the middle and falls
symmetrically, because $C(n,k)=C(n,n-k)$.

For example, row $6$ is $1, 6, 15, 20, 15, 6, 1$, peaking at $k=3=\frac62$. When $n$
is odd the maximum is shared by the two middle entries, as in row $11$, where
$C(11,5)=C(11,6)=462$.
\end{feedback}
\end{prompt}

\begin{prompt}
\textbf{(b)} Claim: the value of $C(n,k)=\dfrac{n!}{k!\,(n-k)!}$ is always a whole
number.

\begin{multipleChoice}
\choice[correct]{True}
\choice{False}
\end{multipleChoice}

\begin{feedback}
True. $C(n,k)$ counts the number of $k$-element subsets of an $n$-element set, and a
count of subsets cannot be fractional. Despite the division in the formula, the
factors of $k!\,(n-k)!$ always cancel completely.
\end{feedback}
\end{prompt}

\begin{prompt}
\textbf{(c)} Claim: the expected value of a random variable must be one of the
possible values the variable can take.

\begin{multipleChoice}
\choice{True}
\choice[correct]{False}
\end{multipleChoice}

\begin{feedback}
False. The expected value is a weighted average, and an average need not be one of
the values being averaged.

The standard example is a single fair die: $E(X)=3.5$, which is not a face of the
die. Likewise, a family might have an expected $2.4$ children.
\end{feedback}
\end{prompt}
\end{problem}

\begin{problem}
Determine which among the following should be counted using permutations and which
should be counted using combinations, and give the count.

\begin{prompt}
\textbf{(a)} A coach wants to select five starters from a team of $14$ players.

\begin{multipleChoice}
\choice{Permutations}
\choice[correct]{Combinations}
\end{multipleChoice}

The count is $\answer{2002}$.

\begin{feedback}
The five starters form an unordered group --- no distinct positions are assigned --- so
\[ C(14,5)=\frac{14!}{5!\,9!}=2002. \]
\end{feedback}
\end{prompt}

\begin{prompt}
\textbf{(b)} From a pack of $20$ distinct Valentine's Day cards, a child specially
selects which cards will go to each of $7$ friends.

\begin{multipleChoice}
\choice[correct]{Permutations}
\choice{Combinations}
\end{multipleChoice}

The count is $\answer{390700800}$.

\begin{hint}
The friends are distinguishable, so it matters not just \emph{which} seven cards are
used but \emph{who gets which}.
\end{hint}

\begin{feedback}
Because the cards are distinct and each friend gets a specific one, this is an
assignment of $7$ cards from $20$ to $7$ labeled recipients --- a permutation:
\[ P(20,7)=20\cdot 19\cdot 18\cdot 17\cdot 16\cdot 15\cdot 14 = 390{,}700{,}800. \]
\end{feedback}
\end{prompt}

\begin{prompt}
\textbf{(c)} Each week, Saturday Night Live comedians propose $40$ sketches, and by
the end of the week they select and order a $10$-sketch lineup.

\begin{multipleChoice}
\choice[correct]{Permutations}
\choice{Combinations}
\end{multipleChoice}

The count is $\answer{3075990524006400}$.

\begin{feedback}
The lineup is ordered, so
\[ P(40,10)=\frac{40!}{30!}=3{,}075{,}990{,}524{,}006{,}400. \]
Compare this with the Exam 2 version of the same situation, where the sketches were
only \emph{selected} and not ordered, giving $C(40,10)=847{,}660{,}528$. Ordering the
lineup multiplies the count by $10!$.
\end{feedback}
\end{prompt}
\end{problem}

\begin{problem}
Complete Pascal's triangle out to the layer that consists of the values of $C(6,k)$
for each $k=0,1,2,3,4,5,6$.

$\answer{1}$, $\answer{6}$, $\answer{15}$, $\answer{20}$, $\answer{15}$,
$\answer{6}$, $\answer{1}$

\begin{hint}
Row $5$ is $1, 5, 10, 10, 5, 1$. Each entry of row $6$ is the sum of the two entries
above it.
\end{hint}

\begin{feedback}
Starting from row $5$, which is $1, 5, 10, 10, 5, 1$, add adjacent pairs:
\[ 1,\ 6,\ 15,\ 20,\ 15,\ 6,\ 1. \]
For instance $C(6,2)=C(5,1)+C(5,2)=5+10=15$.
\end{feedback}
\end{problem}

\begin{problem}
A group of $12$ supervillains in Gotham are organizing a coordinated team of five to
take Batman down.

\begin{prompt}
\textbf{(a)} How many different teams of five supervillains can be selected from the
group of twelve?

$\answer{792}$

\begin{feedback}
\[ C(12,5)=\frac{12!}{5!\,7!}=792. \]
\end{feedback}
\end{prompt}

\begin{prompt}
\textbf{(b)} How many different teams of five are possible if each villain on the
team is given an assigned role (so the same five villains in different roles is
considered distinct)?

$\answer{95040}$

\begin{feedback}
\[ P(12,5)=12\cdot 11\cdot 10\cdot 9\cdot 8 = 95{,}040, \]
which is $5!=120$ times the answer to part (a).
\end{feedback}
\end{prompt}

\begin{prompt}
\textbf{(c)} How many possible teams of five are there if the Joker insists on being
on the team?

$\answer{330}$

\begin{feedback}
The Joker fills one spot, leaving $4$ spots to fill from the other $11$ villains:
\[ C(11,4)=\frac{11!}{4!\,7!}=330. \]
\end{feedback}
\end{prompt}

\begin{prompt}
\textbf{(d)} If the team of five is selected from among the twelve supervillains
entirely at random, what is the probability that the Joker will be selected as part
of the team?

As a percentage rounded to one decimal place: $\answer[tolerance=0.05]{41.7}\%$

\begin{feedback}
Divide part (c) by part (a):
\[
P(\text{Joker})=\frac{C(11,4)}{C(12,5)}=\frac{330}{792}=\frac{5}{12}\approx 0.4167,
\]
about $41.7\%$.

Sanity check: the team is $5$ of the $12$ villains, and by symmetry every villain is
equally likely to be chosen, so the answer had to be $\frac{5}{12}$.
\end{feedback}
\end{prompt}
\end{problem}

\begin{problem}
You've made it to the final round of a trivia game, and all you need to do is answer
the last five questions. These questions are tricky, so suppose your answers tend to
be correct only $40\%$ of the time.

\begin{prompt}
\textbf{(a)} Calculate the probability distribution for the random variable
$X=$ the number of correctly answered questions in this final round. The fifth row of
Pascal's triangle is $(1,5,10,10,5,1)$.

Round each percentage to three decimal places.

$P(X=0)=\answer[tolerance=0.0005]{7.776}\%$ \qquad
$P(X=1)=\answer[tolerance=0.0005]{25.920}\%$ \qquad
$P(X=2)=\answer[tolerance=0.0005]{34.560}\%$

\smallskip
$P(X=3)=\answer[tolerance=0.0005]{23.040}\%$ \qquad
$P(X=4)=\answer[tolerance=0.0005]{7.680}\%$ \qquad
$P(X=5)=\answer[tolerance=0.0005]{1.024}\%$

\begin{hint}
Use $n=5$ and $p=0.4$, so $1-p=0.6$. The six probabilities must total $100\%$.
\end{hint}

\begin{feedback}
With $n=5$, $p=0.4$:
\[
\begin{array}{ll}
P(X=0)=1\,(0.6)^5=0.07776 & \to 7.776\%\\
P(X=1)=5\,(0.4)(0.6)^4=0.25920 & \to 25.920\%\\
P(X=2)=10\,(0.4)^2(0.6)^3=0.34560 & \to 34.560\%\\
P(X=3)=10\,(0.4)^3(0.6)^2=0.23040 & \to 23.040\%\\
P(X=4)=5\,(0.4)^4(0.6)=0.07680 & \to 7.680\%\\
P(X=5)=(0.4)^5=0.01024 & \to 1.024\%
\end{array}
\]
These sum to exactly $100\%$. \checkmark
\end{feedback}
\end{prompt}

\begin{prompt}
\textbf{(b)} What is the probability that you answer \emph{4 or more} of the
questions correctly?

Rounded to three decimal places: $\answer[tolerance=0.0005]{8.704}\%$

\begin{feedback}
\[
P(X\ge 4)=P(X=4)+P(X=5)=0.07680+0.01024=0.08704,
\]
that is, $8.704\%$.
\end{feedback}
\end{prompt}

\begin{prompt}
\textbf{(c)} Suppose a cash prize is awarded based on how many questions in this
round you answer correctly:

\begin{center}
\begin{tabular}{|c|c|c|c|c|c|c|}\hline
$k$ & $0$ & $1$ & $2$ & $3$ & $4$ & $5$ \\\hline
prize for $k$ correct & \$0 & \$50 & \$100 & \$200 & \$500 & \$1000 \\\hline
\end{tabular}
\end{center}

Based on your probability distribution above, what is the expected value for your
prize? Round to two decimal places.

$\$\answer[tolerance=0.005]{142.24}$

\begin{hint}
Multiply each prize by the probability of winning it, then add. Note that you are
taking the expected value of the \emph{prize}, not of the number correct.
\end{hint}

\begin{feedback}
\begin{align*}
E(\text{prize}) &= 0(0.07776)+50(0.25920)+100(0.34560)\\
&\qquad +200(0.23040)+500(0.07680)+1000(0.01024)\\
&= 0 + 12.96 + 34.56 + 46.08 + 38.40 + 10.24\\
&= 142.24.
\end{align*}
So the expected prize is $\$142.24$.

Note this is not a prize you can actually win --- no row of the table pays
$\$142.24$. That is part (c) of Problem 1 in action.
\end{feedback}
\end{prompt}
\end{problem}

\end{document}
